% Generado por sistema/backend ReporteTex.java -- no editar a mano.
% Datos: Base de datos: 148 importado (titulados) / Base de datos: 53 importado (empleadores)
\paragraph{3.2.3 Condiciones de empleo.} \textbf{Qué pide:} que el diseño de la carrera y la calidad de la formación se reflejen en las condiciones de empleo de los graduados; sus indicadores son el tiempo medio para obtener el primer empleo, el destino laboral y la tasa de empleo, y la concordancia entre titulación y empleo. \textbf{Evidencia del estudio:} tasa de empleo de 83.1\% (123 de 148, Sección~\ref{sec:titulados}); tiempo medio para obtener el primer empleo de 2.63 meses, con IC 95\% $[1.85,\ 3.42]$; el 80.6\% de quienes trabajan en una organización lo hace en el sector privado y el rubro más frecuente es ``Desarrollo de Software''; cruces C2 (sin relación significativa, $X_0^2 = 0.002$ frente a $3.841$) y C6 (sin relación significativa, $X_0^2 = 0.227$ frente a $3.841$). El tiempo medio no tiene un valor de referencia en ARCU-SUR, por lo que queda como línea base para el seguimiento.

\textbf{Valoración.} Aspectos favorables: titulados empleados (83.1\%); titulados que consideran pertinentes a su formación las labores de su puesto (90.7\%); titulados que perciben concordancia entre su formación y los requerimientos del mercado (61.9\%). No se identifican aspectos desfavorables.

\paragraph{3.2.1 Resultados.} \textbf{Qué pide:} evaluar el resultado del proceso formativo y utilizar esa evaluación para realizar ajustes. \textbf{Evidencia del estudio:} el 74.1\% de los titulados está de acuerdo con estar satisfecho con la formación recibida (media 2.78 en la escala de 1 a 4); el aspecto a mejorar más señalado es ``Actualizar el plan de estudios'' (71.6\%); cruce C1 (sin relación significativa, $X_0^2 = 0.796$ frente a $5.991$).

\textbf{Valoración.} Aspectos favorables: titulados satisfechos con la formación (74.1\%). No se identifican aspectos desfavorables.

\paragraph{3.2.2 Vinculación y seguimiento a los graduados.} \textbf{Qué pide:} un sistema de seguimiento de los graduados que permita conocer sus condiciones de empleo, e instancias de participación de los graduados para mejorar la carrera. \textbf{Evidencia del estudio:} la encuesta aplicada y la herramienta desarrollada constituyen un mecanismo de seguimiento: cada nueva aplicación se importa y los indicadores se recalculan. El 57.4\% de los titulados manifiesta interés en formar parte de una red oficial de contactos de la Carrera, lo que da una base para esa participación.

\textbf{Valoración.} Aspectos favorables: titulados interesados en la red de contactos (57.4\%). No se identifican aspectos desfavorables.

\paragraph{2.1.2 Perfil de egreso.} \textbf{Qué pide:} coherencia entre el perfil de egreso y la demanda de competencias profesionales expresada por agentes sociales relevantes. \textbf{Evidencia del estudio:} valoración de los empleadores en la escala de perfil profesional (Sección~\ref{sec:empleadores}), donde se considera de acuerdo la suma de ``totalmente'' y ``parcialmente de acuerdo''.

\textbf{Valoración.} Aspectos favorables: ``El título es consistente con las competencias'' (86.8\%); ``Conozco el perfil profesional declarado'' (84.9\%); ``El perfil es coherente con el desempeño'' (79.2\%); ``Evidencian competencias laborales'' (86.8\%); ``Demuestran valores y ética profesional'' (83.0\%). No se identifican aspectos desfavorables.

\paragraph{2.1.7 Actualización curricular.} \textbf{Qué pide:} mecanismos permanentes de actualización del currículo, en consonancia con los desarrollos profesionales. \textbf{Evidencia del estudio:} la asignatura más señalada como ventaja competitiva es ``Base de Datos'' (94.6\%) y la más señalada como poco útil es ``Aplic. Interactivas TV Digital'' (48.6\%); el 71.6\% señala la actualización del plan de estudios como aspecto a mejorar. Las tecnologías y habilidades mencionadas por los empleadores (Sección~\ref{sec:empleadores}) son un insumo directo para esa actualización.

\paragraph{2.4.1 Cursos de actualización profesional y 1.1.5 Posgrado.} \textbf{Qué piden:} extender los conocimientos hacia los graduados mediante cursos de actualización y programas de posgrado afines a la carrera, y promover estudios de posgrado. \textbf{Evidencia del estudio:} el 66.2\% tiene o cursa formación complementaria y el 83.1\% está interesado en un posgrado del área; el área más solicitada es ``IA''; de los interesados, el 41.5\% elegiría la UMSS (EUPG o FCyT); cruces C3 (sin relación significativa, $X_0^2 = 0.769$ frente a $3.841$), C4 (sin relación significativa, $X_0^2 = 1.749$ frente a $3.841$) y C5 (sin relación significativa, $X_0^2 = 2.013$ frente a $9.488$).

\textbf{Valoración.} Aspectos favorables: titulados interesados en un posgrado (83.1\%); titulados que consideran afín la oferta de posgrado de la FCyT (70.9\%). Aspectos desfavorables: interesados que elegirían la UMSS para su posgrado (41.5\%).

\paragraph{2.4.2 Relaciones con el sector público y privado.} \textbf{Qué pide:} que la institución establezca relaciones con empresas y organizaciones públicas y privadas para cooperar en actividades conjuntas. \textbf{Evidencia del estudio:} encuesta a organizaciones de ambos sectores (Sección~\ref{sec:empleadores}).

\textbf{Valoración.} Aspectos favorables: organizaciones que contrataron titulados en los últimos 5 años (79.2\%); ``Consideran nuestra opinión en prácticas y pasantías'' (58.5\%); ``He participado en procesos de retroalimentación'' (52.8\%); ``La Carrera mantiene vínculos que generan oportunidades'' (62.3\%); ``He participado en vinculación entre profesionales'' (50.9\%); ``Recurrimos a la UMSS cuando requerimos profesionales'' (60.4\%). Aspectos desfavorables: ``Las autoridades consultan nuestra opinión'' (47.2\%); ``He participado en vinculación con el entorno laboral'' (47.2\%); ``He participado en difusión de oportunidades laborales'' (47.2\%); ``He sido consultado sobre el desempeño de los titulados'' (47.2\%); ``Han consultado nuestra opinión sobre la pertinencia'' (43.4\%); ``He sido consultado en la autoevaluación de la Carrera'' (35.8\%).

\paragraph{1.2.2 Sistemas de información para la toma de decisiones.} \textbf{Qué pide:} sistemas con información relevante, confiable y actualizada que respalden las decisiones institucionales. \textbf{Evidencia del estudio:} la herramienta desarrollada importa las respuestas, elimina los datos personales, calcula los indicadores y cruces de este informe y exporta los resultados, por lo que puede usarse de forma permanente para actualizar esta información.

